% Same-history physical stress transfer and Scribe assembly.
\section{The same-history physical stress transfer}
\label{sec:ex-lattice-trace-scribe-producer}

This section proves the physical stress and flux passage used by the
selected trace gallery from the explicit response, sensor, fourth-order,
and weighted trace receipts displayed below. Those paid receipts are
producer inputs; the arguments here show how they combine. All limits below use one selected history, one
ultrafilter, one hierarchy, and one scaling. In particular, the physical
stress, the correction tensors, and the comparison observables are never
reselected on separate subsequences. The quantitative estimates entering the
construction include the paid response and comparison estimates from
Sections~\ref{sec:ex-lattice-trace-weighted-estimates},
\ref{sec:ex-lattice-trace-crown-buffer}, and
\ref{sec:ex-lattice-trace-eve-producer}. The additional sensor and branch
receipts are stated explicitly below. Given these inputs, this section proves
how the estimates produce the conserved tensor sequence and transfer its
physical normalization to the selected table.

Write $\kappa>0$ for the slope constant and let $\mathcal O$ be the
distinguished unit field in the reference grade-four stress table. The
physical coefficient is $d_2=-\kappa/4$.

\subsection{The two-direction conserved correction}

Let $E$ be the finite-dimensional real space of centered, signed-equivariant,
spin-symmetric grade-four tensors for the selected table (before imposing
the divergence equation), let $V$ be the finite-dimensional
response space, and let $\mathcal B$ be the grade-five profile space. Let
$D:E\to V$ be the divergence map. Under the reference-rank and optical-comparison receipts
recorded in Section~\ref{sec:ex-lattice-trace-producer}, the finite-dimensional
argument there gives $\dim\ker D=1$. We use that conclusion here. The
interacting divergence remains fixed; the comparison varies the optical map,
and its positive scale and invertible chart preserve the kernel.

The construction solves the correction equations at each
large index, rather than taking an unspecified kernel vector. Let
$\mathsf K_i:E\to V$ be the selected comparison map, let
$s_i:\mathcal B\to V$ be the paid response map, and let
$L_i,\Omega_i,\Delta_i$ be the
principal and two correction directions. The correction map is
\[
 \mathsf S_i(c_0,c_1)=c_0\Omega_i+c_1\Delta_i,
 \qquad \mathsf B_i=\mathsf K_i\mathsf S_i.
\]
Let $q_i>0$ be the selected effective coupling and $\rho_i$ its positive
response-size price; these are the scales in the paid bounds below.
For $a=0,1$, the corrected response estimates compare the corresponding
column with the same-history linear response on a paid forecast
$w_i^a\in\mathcal B$:
\[
 \|\mathsf B_i e_a-s_i(w_i^a)\|\longrightarrow0,
 \qquad w_i^a\longrightarrow v_a,
 \qquad s_i\longrightarrow\mathsf V.
\]
The profiles $v_0,v_1$ are the two free correction directions, and
$\mathsf W(c_0,c_1)=c_0v_0+c_1v_1$. Hence
\[
 \mathsf B_i\longrightarrow
 \mathsf B_\infty=\mathsf V\mathsf W,
 \qquad
 \mathsf V\mathsf W(c_0,c_1)=c_0\mathsf V(v_0)+c_1\mathsf V(v_1),
\]
where $\mathsf V$ is the free vector-response map and $\mathsf W$ is the
two-column reference correction. Indeed, for each $a$, the displayed
response estimate and continuity of the limiting response give
$\mathsf B_i e_a\to\mathsf V(v_a)$. Since the domain has basis
$e_0,e_1$, convergence of both columns implies operator-norm convergence.
The free classification proves
$\mathsf B_\infty$ is an isomorphism. In finite dimension, if
$\|\mathsf B_i-\mathsf B_\infty\|<1/(2\|\mathsf B_\infty^{-1}\|)$,
then for every $c$,
\[
 \|\mathsf B_i c\|\ge
 \|\mathsf B_\infty^{-1}\|^{-1}\|c\|-\|
   (\mathsf B_i-\mathsf B_\infty)c\|
 \ge \frac{\|c\|}{2\|\mathsf B_\infty^{-1}\|}.
\]
Thus $\mathsf B_i$ is injective and hence bijective, and its inverse norms
are uniformly bounded on a tail. Let $X_i=q_i\mathsf K_i(L_i)$ be the
uncorrected forcing. The paid response-size estimate gives
$\|X_i\|\le C\rho_i$, where $\rho_i$ is the positive response-size
price. Define
\[
 c_i=-\mathsf B_i^{-1}X_i,\qquad
 T_i=q_iL_i+\mathsf S_i(c_i).
\]
Then $\mathsf K_i(T_i)=0$, and the optical comparison identity gives
$DT_i=0$: on the selected tail, $\mathsf K_i(Q)=0$ is equivalent to
$DQ=0$ for every $Q\in E$. The inverse bound gives
$\|c_i\|\le C'\rho_i$.

The scalar calibration uses the same paid functional $\mu_i\in V^*$ that controls
the response price. Let $\Lambda\in\mathcal B^*$ be the free Ward functional.
At chart level $k_i$, write $t_i$ for its positive sensor
scale and $z_i=\widehat\beta_{k_i}$ for its clock weight. The price comparison
gives $t_i\to\infty$, $\rho_i\le t_i\le z_i$. The sensor inputs
include a constant $C_\mu$ with $\|\mu_i\|\le C_\mu$ eventually, and, for
every $\delta>0$, the eventual estimate
\[
 t_i\left|\mu_i(s_i y)-\Lambda(y)\right|
       \le\delta\|y\|\quad\text{for all }y\in\mathcal B.
\]
We spell out how these inputs
give the three limits needed for calibration. Let $Y_i$ be the principal
forecast. The paid response comparison, penumbra estimate, and principal
Ward limit are
\begin{equation}
 q_i\mathsf K_i(L_i)-s_i(Y_i)\longrightarrow0,\qquad
 \|Y_i\|\le C\rho_i\le Ct_i,\qquad
 \Lambda(Y_i)\longrightarrow-\kappa.
 \label{eq:ex-lattice-scribe-principal-response-input}
\end{equation}
The first convergence and the uniform bound on $\mu_i$ imply
$\mu_i(X_i)-\mu_i(s_i(Y_i))\to0$. For any $\delta>0$, the sensor estimate
then gives eventually
\[
 |\mu_i(s_i(Y_i))-\Lambda(Y_i)|
       \le \delta\|Y_i\|/t_i\le C\delta.
\]
Taking the limit superior and then $\delta\downarrow0$ proves
$\mu_i(X_i)\to-\kappa$; no norm convergence of the principal forecast is
needed.

For a correction branch $P\in\{\Omega,\Delta\}$, let $U_i^P\in\mathcal B$
be its finite-volume profile and $v_P\in\mathcal B$ its reference profile.
The comparison and
Ward estimates needed are
\begin{equation}
 U_i^P\to v_P,\qquad
 z_i\|\mathsf K_i(P_i)-s_i(U_i^P)\|\to0,\qquad
 t_i|\Lambda(U_i^P)-\Lambda(v_P)|\to0,
 \label{eq:ex-lattice-scribe-correction-response-input}
\end{equation}
where $P_i$ denotes the corresponding correction tensor. Since $U_i^P$ is
bounded, the uniform sensor estimate implies
$t_i|\mu_i(s_i(U_i^P))-\Lambda(U_i^P)|\to0$. Using $t_i\le z_i$ and the
uniform operator bound for $\mu_i$ in the middle estimate of
(\ref{eq:ex-lattice-scribe-correction-response-input}) gives
\[
 t_i\left|\mu_i(\mathsf K_i(P_i))-\Lambda(v_P)\right|\longrightarrow0.
\]
The reference Ward values are $\Lambda(v_\Omega)=0$ and
$\Lambda(v_\Delta)=-4$. Thus the principal and two correction limits are
now established from the explicitly stated paid inputs:
\begin{equation}
 \mu_i(X_i)\to-\kappa,\qquad
 t_i\mu_i(\mathsf K_i\Omega_i)\to0,\qquad
 t_i\bigl(\mu_i(\mathsf K_i\Delta_i)+4\bigr)\to0.
 \label{eq:ex-lattice-scribe-sensor-limits}
\end{equation}
Applying $\mu_i$ to the exact cancellation $X_i+\mathsf B_i c_i=0$ gives
\[
 \mu_i(X_i)+c_i^0\mu_i(\mathsf K_i\Omega_i)
       +c_i^1\mu_i(\mathsf K_i\Delta_i)=0.
\]
Set $\alpha_i=t_i\mu_i(\mathsf K_i\Omega_i)$ and
$\beta_i=t_i(\mu_i(\mathsf K_i\Delta_i)+4)$. Since
$\|c_i\|\le C\rho_i\le Ct_i$, both $c_i^0/t_i$ and $c_i^1/t_i$ are
bounded. The exact equation rewrites as
\[
 0=\mu_i(X_i)-4c_i^1+
       \frac{c_i^0}{t_i}\alpha_i+\frac{c_i^1}{t_i}\beta_i.
\]
The last two terms tend to zero by
(\ref{eq:ex-lattice-scribe-sensor-limits}), so
$\mu_i(X_i)\to-\kappa$ implies
\[
 c_i^1\longrightarrow-\frac{\kappa}{4}.
\]
This is the stable-inversion and coefficient-extraction argument of
Section~\ref{sec:ex-lattice-trace-producer}.

The source constructs the correction using the same three paid tensor
families at every index: a principal tensor $L_i$ and two correction tensors
$\Omega_i,\Delta_i$. Let
$\mathcal P_i$ be the scaled pullback of a stress tensor, and put
$\mathcal Y_i=q_i^{-1}\mathcal P_i$. The scaled pullback is the
chart map multiplied by its positive crown normalization; by definition,
$\mathcal P_i=q_i\mathcal Y_i$ on every tensor. In the source,
$q_i$ is the effective coupling at the selected level and the crown is
$q_i$ times the fourth power of the chart radius. The defining corrected tensor is
\begin{equation}
 T_i=q_iL_i+c_i^0\Omega_i+c_i^1\Delta_i.
 \label{eq:ex-lattice-scribe-corrected-tensor}
\end{equation}
For the three tensor sequences $Q_i\in\{L_i,\Omega_i,\Delta_i\}$, let
their reference branches be $\nu_L=0$, $\nu_\Omega=-$, and
$\nu_\Delta=+$. Let $W_i^Q(a,b)$ be the ordered grade-four label for the
two tensor indices $a,b$ obtained by the finite-volume response of the
corresponding sequence. These ordered
labels are the entries of a finite array; no spatial sum is hidden in this
notation. Let $\mathsf M_i$ be the finite-volume fourth-order coefficient
map, obtained by applying the centered chart pullback to a grade-four label.
It is a continuous linear map
$\mathsf M_i:\mathcal L_4\to\mathcal F_4$, not a scalar moment matrix.
Each $W_i^Q(a,b)$ belongs to $\mathcal L_4$. For every sufficiently large
$i$ (so the selected level exceeds the fixed reference cutoff), symmetrizing
the two ordered indices gives the chart formula
\begin{equation}
 [\mathcal P_i(Q_i)]_{ab}=
 \frac12\left(\mathsf M_i(W_i^Q(a,b))+\mathsf M_i(W_i^Q(b,a))\right).
 \label{eq:ex-lattice-scribe-phase-formula}
\end{equation}
The selected coefficient convergence gives $W_i^Q(a,b)\to
W_{\nu_Q}(a,b)$, and the fourth-order maps converge in operator norm,
$\mathsf M_i\to\mathsf M_\infty$. Thus, entry by entry in the finite
stress array,
\[
 \mathcal P_i(Q_i)\longrightarrow\mathcal B_{\nu_Q},\qquad
 (\mathcal B_\nu)_{ab}:=
 \frac12\left(\mathsf M_\infty(W_\nu(a,b))+
                \mathsf M_\infty(W_\nu(b,a))\right).
\]
This follows directly from continuity of the finite bilinear expression in
(\ref{eq:ex-lattice-scribe-phase-formula}); no limit is interchanged with
an infinite sum. The explicit free tensor calculation identifies these
arrays with the reference material tensors. Write
\[
 M_0=\mathcal B_0,\qquad M_\Omega=\mathcal B_-,\qquad M_1=\mathcal B_+.
\]
The false-branch tensor need not vanish; it is its coefficient
$c_i^0/q_i\to0$ that removes it from the normalized field limit. The three
branch limits are on the same history.

These are the field receipts used below:
\begin{align}
 \mathcal P_i(L_i)&\longrightarrow M_0,&
 \mathcal P_i(\Omega_i)&\longrightarrow M_\Omega,&
 \mathcal P_i(\Delta_i)&\longrightarrow M_1,
 &\frac{\rho_i}{q_i}&\longrightarrow0,
 \label{eq:ex-lattice-scribe-field-receipts}
\end{align}
The displayed ratio is the selected-scale washout input: the response price
is bounded by the one-sixth subpower clock price, whose ratio to the effective
coupling tends to zero along this scale. Since $M_0$ is the prescribed
material tensor and $|c_i^0|+|c_i^1|\le C\rho_i$, it follows that
$c_i^a/q_i\to0$ for $a=0,1$. Linearity gives
the exact identity
\[
 \mathcal Y_i(T_i)=\mathcal P_i(L_i)
   +\frac{c_i^0}{q_i}\mathcal P_i(\Omega_i)
   +\frac{c_i^1}{q_i}\mathcal P_i(\Delta_i).
\]
The principal, false-correction, and true-correction rescaled tensors converge
to $M_0,M_\Omega,M_1$, respectively.
In particular the two correction phase tensors are bounded, and the exact
identity proves
\begin{equation}
 \mathcal Y_i(T_i)\longrightarrow M_0.
 \label{eq:ex-lattice-scribe-field-limit}
\end{equation}
For the scalar trace, the true correction profile is the true material
correction, whose diagonal trace is $4\mathcal O$ by the finite tensor
identity used in the source: the diagonal sum is
$4(0+4\cdot\frac14)\mathcal O=4\mathcal O$. The source bounds
$q_i\le2z_i$ and $\rho_i\le t_i\le z_i$ and supplies
\[
 z_i\operatorname{tr}\mathcal P_i(L_i)\to0,\quad
 z_i\operatorname{tr}\mathcal P_i(\Omega_i)\to0,\quad
 |c_i^0|\le C\rho_i,\quad
 \operatorname{tr}\mathcal P_i(\Delta_i)\to4\mathcal O,
 \quad c_i^1\to-\kappa/4.
\]
To see exactly how the weighted estimates enter, use
$\mathcal P_i(T_i)=q_i\mathcal P_i(L_i)
+c_i^0\mathcal P_i(\Omega_i)
+c_i^1\mathcal P_i(\Delta_i)$. The principal contribution tends
to zero since $q_i\le2z_i$ and the exponent-one principal trace estimate
gives $z_i\operatorname{tr}\mathcal P_i(L_i)\to0$. For the false
correction, $|c_i^0|\le C\rho_i\le Cz_i$ and the exponent-one
weighted false-sector estimate gives
$z_i\operatorname{tr}\mathcal P_i(\Omega_i)\to0$.
Finally $c_i^1\to-\kappa/4$ and
$\operatorname{tr}\mathcal P_i(\Delta_i)\to4\mathcal O$.
Thus applying the trace to (\ref{eq:ex-lattice-scribe-corrected-tensor})
and using these three contributions gives
\begin{equation}
 \operatorname{tr}\mathcal P_i(T_i)\longrightarrow-\kappa\mathcal O.
 \label{eq:ex-lattice-scribe-flux-limit}
\end{equation}
Indeed, the principal term tends to zero, the first correction term is
bounded by $C z_i|\operatorname{tr}\mathcal P_i(\Omega_i)|$, and the last
term tends to $(-\kappa/4)(4\mathcal O)$. This establishes the physical
field and flux limits from the exact three-term correction, with no
interchange of limits or choice of history.

\subsection{Proportionality transfers the physical normalization}

We now show that the corrected physical tensor and the selected table
stress have the same limiting normalization. Let $A$ be the selected
algebra, and let $R$ be its centered table stress
$R=\operatorname{reCenter}_S(A.\operatorname{energy})\in E$, whose entries
are conserved and whose flux is the physical Ward flux. This is the stress
used in the physical-shell theorem. Let $g_i$ be the explicit conserved
correction sequence just constructed, viewed in the same one-dimensional
kernel of $D$. Assume its normalized field
and flux limits are (\ref{eq:ex-lattice-scribe-field-limit}) and
(\ref{eq:ex-lattice-scribe-flux-limit}), and that the table stress has the
same nonzero shell-flux limit. We prove the transfer, spelling out the
nonzero test and the scalar ratio.

The reference observer-separation calculation supplies an axis $a$, an
observer $o$, and a reference grade-four field direction $v$ such that
$y=o.test_5(D_a v)\ne0$, where $D_a$ is the reference derivative in direction
$a$. Let $\chi$ be the cut and $x$ the base point belonging to this same
observer. Its defining test identity identifies this number with
$\operatorname{shellGrad}(M_0,\chi,a,x,v)$. We keep this exact observer and
cut for both stress sequences; no new test or window is selected in the
limit argument.

For a stress tensor $Q$, write the exact selected shell measurement as
\[
 \mathscr F_i(Q)=r_i\operatorname{shellGrad}_S\!\left(
 Q,\operatorname{stretch}(r_i,\chi),a,r_i x,
 \operatorname{pins}_i\circ q,U_i v\right),
\]
where $r_i$ is the chart radius, $q$ is the observer's fixed test window,
and $U_i$ is the selected grade-four chart isomorphism. The selected
chart's tested convergence and the comparison-of-jets estimate give, for
the physical tensor and for $g_i$, respectively,
\begin{equation}
 f_i:=\mathscr F_i(R)\to y,
 \qquad \mathscr F_i(g_i)\to y.
 \label{eq:ex-lattice-scribe-common-test}
\end{equation}
The argument uses the same selected scale and the same observer test in both
limits. The physical flux excludes $R=0$: if $R$ vanished, linearity of
shell differentiation would force $f_i=0$ at every index, contrary to
$y\ne0$.

Since $\ker D$ is one-dimensional, $g_i=\lambda_iR$ for all sufficiently
large $i$. The second limit in (\ref{eq:ex-lattice-scribe-common-test})
therefore reads $\lambda_i f_i\to y$, while the first reads $f_i\to y$.
As $y\ne0$, choose a scalar coordinate functional $\ell$ with $\ell(y)\ne0$.
Then
\[
 \lambda_i=\frac{\ell(\lambda_i f_i)}{\ell(f_i)}\longrightarrow
 \frac{\ell(y)}{\ell(y)}=1.
\]
Thus $g_i=\lambda_iR$ and $\lambda_i\to1$. On a late tail
$\lambda_i\ne0$ and $\lambda_i^{-1}\to1$. Since $\mathcal Y_i$
and the weighted trace are linear,
\[
 \mathcal Y_i(R)=\lambda_i^{-1}\mathcal Y_i(g_i),
 \qquad
 \operatorname{tr}\mathcal P_i(R)=\lambda_i^{-1}
       \operatorname{tr}\mathcal P_i(g_i)
\]
on that tail. This proves the physical field and flux limits. In particular,
\begin{equation}
 \mathcal Y_i(R)\longrightarrow M_0,
 \qquad \operatorname{tr}\mathcal P_i(R)\longrightarrow-\kappa\mathcal O.
 \label{eq:ex-lattice-scribe-physical-transfer}
\end{equation}
This is the physical-array step used by the chart gallery. The dimension-one
kernel alone would not fix the scalar; the common nonzero shell test fixes it.

\subsection{Comparison tables and the full gallery record}

Let $T$ be any comparison table that respects the selected table, and let
$\Theta_T$ be its trace observable on $A.\operatorname{energy}$. Its exact
centered decomposition is
\[
 \Theta_T=\operatorname{tr}(\operatorname{reCenter}_S(A.\operatorname{energy}))-c_T\,1,
 \qquad c_T\in\mathbb C.
\]
The physical-transfer limit (\ref{eq:ex-lattice-scribe-physical-transfer})
applies to the first term. The scalar-unit term vanishes after zooming:
the selected scale satisfies
\[
 q_i r_i^4\,1\longrightarrow0,
 \qquad q_i>0,
\]
and the selected chart eventually identifies its unit with the reference
unit. The zoom maps are linear, so the same vanishing holds after
multiplication by $c_T$. Therefore the zoomed comparison observable converges to the same
prescribed multiple of the unit as the physical stress. This argument is
uniform in the algebraic identity for $T$; only the two scalar sequences
used to represent its component in the distinguished two-dimensional
stress sector depend on $T$.

More explicitly, the scalar representation is
\[
 \iota(e,d)=e\,1+4d\,\mathcal O.
\]
The finite reference calculation shows that $1$ and $\mathcal O$ are
linearly independent, so $\iota$ is injective and has a continuous left
inverse on its range.
The limiting comparison tensor is $4d_2\mathcal O$, where
$d_2=-\kappa/4$, so its coordinates are $(0,d_2)$. Put
$Z_i=q_i r_i^4(J.U_i)^{-1}\Theta_T$. The selected comparison receipt gives
$Z_i\in\operatorname{range}(\iota)$ eventually, separately from the limit
$Z_i\to4d_2\mathcal O$. On that tail let $(a_i,b_i)=\iota^{-1}(Z_i)$;
continuity of the fixed left inverse gives $a_i\to0$ and $b_i\to d_2$.
Define $e_i=a_i/q_i$ and $d_i=b_i/q_i$. Thus, eventually,
\[
 r_i^4(J.U_i)^{-1}\Theta_T=\iota(e_i,d_i),
 \qquad q_i e_i\to0,\qquad q_i d_i\to d_2.
\]
The definition of the zoom then gives
$q_i r_i^4(J.U_i)^{-1}\Theta_T\to4d_2\mathcal O$.

The exact record contains $q_i\to\infty$ and
$\log(\operatorname{Zoom}_i)/q_i\to\kappa^{-1}$, where
$\operatorname{Zoom}_i$ is the positive zoom scale of the selected chart;
the same finite-volume
chart and hierarchy; the physical stress limit; and, for every comparison
table, the late scalar representation and two coefficient limits just
proved. The fixed-scale case uses the selected
chart directly. For an arbitrary positive null sequence, the paid scale
selection first chooses a strictly increasing subsequence and a scaling
whose radii are that subsequence; the same history is then passed to the
finite-volume chart producer, and the physical transfer above applies
without another extraction. Thus, under the same paid response, sensor,
branch, and scale-selection inputs, the output of the selected-run producer
is the full Scribe record, rather than only its finite-dimensional chart
part.
